% Part: incompleteness
% Chapter: representability-in-q
% Section: c

\documentclass[../../../include/open-logic-section]{subfiles}

\begin{document}

\olfileid[pt]{inc}{req}{cee}
\olsection{As funções $C$}

Seja $C$ o menor conjunto de funções que contém
\begin{enumerate}
\item $0$,
\item a função sucessor,
\item a adição,
\item a multiplicação,
\item as projeções e
\item a função característica da igualdade, $\Char{=}$;
\end{enumerate}
e que é fechado sob
\begin{enumerate}
\item composição e
\item busca ilimitada, aplicada a funções
regulares.
\end{enumerate}
Lembre que essa última restrição significa simplesmente que só se pode
usar a operação $\mu$ quando o resultado é total. Compare isso com a
definição das funções \emph{recursivas gerais}: aqui acrescentamos a
adição, a multiplicação e $\Char{=}$, mas retiramos a recursão
primitiva.

É claro que toda função em $C$ é recursiva, pois a adição, a
multiplicação e $\Char{=}$ o são. Mostraremos que a recíproca também é
verdadeira; isso equivale a afirmar que, com os demais recursos de $C$,
podemos efetuar a recursão primitiva.

\end{document}
